% ============================================================================
% Fichier  : partie2/P02_C04_gouvernance.tex
% Titre    : La Gouvernance de la Sécurité
% Auteur   : MINKA MI NGUIDJOÏ Thierry Emmanuel
% ============================================================================
% Sources : ISO/IEC 27001:2022 (clause 5) ; COBIT 2019. Mairie d'Eséka : fictif.
% ============================================================================

\chapter{La Gouvernance de la Sécurité}
\label{chap:gouvernance}
\chaptermark{La gouvernance}

\epigraph{%
    « Sans responsable nommé, la sécurité est l'affaire de personne. »%
}{Principe du cours}

\niveauChapitre{Fondamental}{La partie I.}

\begin{boxobjectifs}
\begin{itemize}
    \item Définir la gouvernance de la sécurité et la distinguer du management
          opérationnel.
    \item Préciser le rôle du \sigle{RSSI}, de la direction et du comité de
          sécurité.
    \item Comprendre le rattachement et l'indépendance du \sigle{RSSI}.
    \item Établir une matrice de responsabilités (\sigle{RACI}).
    \item Relier la gouvernance à l'engagement de la direction exigé par la
          norme.
\end{itemize}
\end{boxobjectifs}

\bigskip

% ============================================================================
\section*{Question directrice : qui décide, qui exécute, qui contrôle ?}
% ============================================================================

Nommer une \sigle{RSSI} ne suffit pas~: encore faut-il définir ses pouvoirs,
son rattachement et ses relations avec la direction et les services. Ce
chapitre organise la gouvernance de la sécurité à la mairie.

\medskip
\textbf{Plan du chapitre.}
\begin{enumerate}
    \item Gouverner et manager (section~\ref{sec:c04-gouverner}).
    \item Le RSSI (section~\ref{sec:c04-rssi}).
    \item La direction et le comité de sécurité (section~\ref{sec:c04-comite}).
    \item La matrice de responsabilités (section~\ref{sec:c04-raci}).
    \item L'engagement de la direction (section~\ref{sec:c04-engagement}).
\end{enumerate}

\separateur

% ============================================================================
\section{Gouverner et Manager}
\label{sec:c04-gouverner}
% ============================================================================

\begin{boxdef}
\termecle{Gouvernance de la sécurité}~: ensemble des responsabilités et
pratiques par lesquelles la direction oriente la sécurité, fixe les objectifs
et s'assure que les risques sont maîtrisés. Elle décide~; le management
exécute.
\end{boxdef}

La gouvernance répond à « quoi et pourquoi »~; le management opérationnel à
« comment ». Confondre les deux conduit soit à une direction qui bricole de la
technique, soit à une technique sans mandat. La norme sépare nettement le rôle
d'orientation (direction) du rôle de pilotage (\sigle{RSSI}).

% ============================================================================
\section{Le RSSI}
\label{sec:c04-rssi}
% ============================================================================

\begin{tableaucours}{L{3.0cm}L{9.2cm}}{%
    \tableauHeader{Mission} &
    \tableauHeader{Contenu}
}
Définir & Proposer la stratégie et la politique de sécurité \\
Piloter & Conduire l'appréciation des risques et le plan de traitement \\
Coordonner & Mettre en relation direction, DSI, agents, prestataires \\
Contrôler & Vérifier le respect des règles, suivre les indicateurs \\
Conseiller & Éclairer la direction pour ses arbitrages \\
\end{tableaucours}
\vspace{-0.8em}
\begin{center}{\small\itshape Tableau~4.1. Les missions du RSSI}\end{center}

\begin{boxprocedure}
\textbf{Constatation.} À la mairie, Mme~\textsc{Abena} (RSSI) est aussi chef du
service informatique, sous l'autorité du \sigle{DSI}.

\medskip
\textbf{Analyse} (confiance moyenne). Ce rattachement crée un conflit~: juger
la sécurité des systèmes que l'on exploite soi-même, c'est être juge et
partie. Il est préférable que la \sigle{RSSI} rende compte directement au
maire, pour garder son indépendance de contrôle.
\end{boxprocedure}

% ============================================================================
\section{La Direction et le Comité de Sécurité}
\label{sec:c04-comite}
% ============================================================================

La direction (le maire) impulse, valide la politique, alloue les ressources et
assume la responsabilité finale. Pour instruire ses décisions, elle s'appuie
sur un \termecle{comité de sécurité}~: instance périodique qui réunit maire,
\sigle{RSSI}, \sigle{DSI} et responsables de services, arbitre les risques et
suit le plan d'action.

\begin{boxmetier}{Direction générale}
\textbf{Présider, pas assister.} Un comité de sécurité que le maire délègue
systématiquement perd son autorité. La présence de la direction signale que la
sécurité est une priorité, non une affaire de techniciens.
\end{boxmetier}

% ============================================================================
\section{La Matrice de Responsabilités}
\label{sec:c04-raci}
% ============================================================================

\begin{boxdef}
Une matrice \termecle{RACI} attribue, pour chaque activité, qui est
\emph{Responsable} (fait), \emph{Approbateur} (décide), \emph{Consulté} et
\emph{Informé}. Elle supprime les zones où « on croyait que l'autre le
faisait ».
\end{boxdef}

\begin{tableaucours}{L{4.0cm}L{2.0cm}L{2.0cm}L{2.0cm}}{%
    \tableauHeader{Activité} &
    \tableauHeader{Maire} &
    \tableauHeader{RSSI} &
    \tableauHeader{DSI}
}
Valider la politique & A & R & C \\
Apprécier les risques & I & R & C \\
Déployer une mesure & I & A & R \\
Suivre les indicateurs & I & R & C \\
\end{tableaucours}
\vspace{-0.8em}
\begin{center}{\small\itshape Tableau~4.2. Extrait de la matrice RACI de la mairie}\end{center}

% ============================================================================
\section{L'Engagement de la Direction}
\label{sec:c04-engagement}
% ============================================================================

La gouvernance n'existe que si la direction s'engage réellement. La clause~5
de \norme{ISO/IEC 27001} l'exige~: la direction doit démontrer son
\termecle{leadership} en établissant la politique, en attribuant les rôles et
en fournissant les ressources (Chapitre~\ref{chap:clauses46}).

\begin{boxjuris}
En droit camerounais comme dans la norme, la responsabilité de la sécurité des
données remonte au responsable de traitement. Le maire ne peut se retrancher
derrière ses services~: la gouvernance formalise cette responsabilité, elle ne
la dilue pas.
\end{boxjuris}

\separateur

\begin{boxresume}
\textbf{Ce qu'il faut retenir de ce chapitre~:}
\begin{itemize}
    \item La \textbf{gouvernance} oriente et décide~; le management exécute.
    \item Le \textbf{RSSI} définit, pilote, coordonne, contrôle, conseille~;
          son indépendance suppose un bon rattachement.
    \item La \textbf{direction} impulse et assume~; le \textbf{comité de
          sécurité} instruit ses décisions.
    \item La \textbf{matrice RACI} supprime les responsabilités floues.
    \item La norme exige l'\textbf{engagement} de la direction (clause~5).
\end{itemize}
\end{boxresume}

\begin{boxinfo}
\textbf{Points de contrôle de l'élève.} Je sais~: distinguer gouverner et
manager~; énoncer les missions du RSSI~; expliquer le rôle du comité~; lire une
matrice RACI~; relier la gouvernance à la clause~5.
\end{boxinfo}

\bigskip

\begin{boxexo}[title={\ding{45}\quad Exercice~4.1. Gouverner ou manager \niveaubase}]
Classez~: fixer la politique~; configurer un pare-feu~; arbitrer un budget~;
appliquer un correctif. Gouvernance ou management~?
\end{boxexo}

\begin{boxexo}[title={\ding{45}\quad Exercice~4.2. Missions du RSSI \niveaubase}]
Citez cinq missions du \sigle{RSSI} et donnez un exemple pour l'une d'elles.
\end{boxexo}

\begin{boxexo}[title={\ding{45}\quad Exercice~4.3. Le comité \niveaubase}]
Qui siège au comité de sécurité de la mairie, et à quoi sert-il~?
\end{boxexo}

\begin{boxexo}[title={\ding{45}\quad Exercice~4.4. Juge et partie \niveaumoyen}]
Pourquoi rattacher la RSSI au DSI pose-t-il un problème d'indépendance~?
Proposez un meilleur rattachement.
\end{boxexo}

\begin{boxexo}[title={\ding{45}\quad Exercice~4.5. RACI \niveaumoyen}]
Complétez une ligne RACI pour l'activité « gérer les accès des agents ».
\end{boxexo}

\begin{boxexo}[title={\ding{45}\quad Exercice~4.6. Leadership \niveaumoyen}]
Donnez trois actes concrets par lesquels le maire démontre son engagement.
\end{boxexo}

\begin{boxexo}[title={\ding{45}\quad Exercice~4.7. Délégation \niveaumoyen}]
Que peut déléguer le maire, que ne peut-il pas déléguer~? Justifiez.
\end{boxexo}

\begin{boxexo}[title={\ding{45}\quad Exercice~4.8. Zone floue \niveaumoyen}]
Donnez un exemple de responsabilité floue à la mairie et montrez comment la
RACI la lève.
\end{boxexo}

\begin{boxexo}[title={\ding{45}\quad Exercice~4.9. Charte de comité \niveauexpert}]
Rédigez la charte du comité de sécurité de la mairie (composition, fréquence,
missions, décisions).
\end{boxexo}

\begin{boxexo}[title={\ding{45}\quad Exercice~4.10. Fiche de poste \niveauexpert}]
Rédigez la fiche de poste de la RSSI de la mairie d'Eséka, avec rattachement et
pouvoirs.
\end{boxexo}

\bigskip

\begin{boxinfo}
\textbf{Pour aller plus loin.}
\begin{itemize}
    \item ISO/IEC 27001, clause~5~\cite{iso27001}~; COBIT 2019~\cite{cobit2019}.
\end{itemize}
\end{boxinfo}

\begin{boxinfo}
\textbf{Transition.} La gouvernance s'exprime d'abord dans un document~: le
Chapitre~\ref{chap:politique} construit la politique de sécurité.
\end{boxinfo}
